\section{Level III: PathologyFL}
\begin{equation}
\begin{aligned}
\Delta\theta_i^{(t)}
&= \operatorname{LocalTrain}\!\left(\theta^{(t)},\mathcal{D}_i;E_i,\mu_i,\varepsilon_i,\delta_i\right), \\
a_i^{(t)}
&= \mathbf{1}\!\left[
\operatorname{finite}(\Delta\theta_i^{(t)})
\land \tau_i^{(t)}\le T_i^{(t)}
\land \neg\operatorname{Byz}_i^{(t)}
\land \operatorname{Integrity}_i^{(t)}
\right], \\
\widetilde w_i^{(t)}
&= a_i^{(t)}w_i^{(t)}e^{-\gamma\tau_i^{(t)}},
&
\bar w_i^{(t)}
&= \frac{\widetilde w_i^{(t)}}{\sum_{j\in\mathcal{I}_t}\widetilde w_j^{(t)}}, \\
\theta^{(t+1)}
&= \theta^{(t)}
+ \operatorname{SecureAgg}\!\left(
\left\{\bar w_i^{(t)}\Delta\theta_i^{(t)}\right\}_{i\in\mathcal{I}_t}
\right), \\
H_t
&= \operatorname{SHA256}\!\left(
H_{t-1}\Vert t\Vert\mathcal{I}_t\Vert
\{\operatorname{SHA256}(\Delta\theta_i^{(t)})\}_{i\in\mathcal{I}_t}
\Vert\operatorname{SHA256}(\theta^{(t+1)})
\right).
\end{aligned}
\label{eq:pathologyfl}
\end{equation}


The third aggregation level is institutional aggregation: how multi-institution
updates, validation priorities, contribution signals, and safety constraints are
aggregated into a collaborative model. The program's contribution here is
PathologyFL, a pathology-specific federated-learning research framework, and
FAIR-WEIGHTS-H, an auditable hybrid institutional-weighting protocol.

\subsection{Architecture}

PathologyFL implements a coordinator-client architecture under
\texttt{src/features/federated/pathology\_fl/}:

\begin{itemize}
\item \textbf{Coordinator:} a production orchestrator with round lifecycle,
  client selection, Byzantine filtering, a tamper-evident audit-log hash chain,
  and versioned checkpoints; a failure handler with timeouts, retries,
  blacklisting, and recovery; monitoring (Prometheus, TensorBoard, convergence
  detection, alerts); and a model registry with SHA-256 integrity and rollback.
\item \textbf{Client:} a production hospital client with local training,
  differential-privacy integration \cite{abadi2016dp}, gRPC secure
  communication, and round
  participation; a local trainer; a resource manager; and a PACS connector.
\item \textbf{Aggregator:} FedAvg \cite{mcmahan2017fedavg},
  FedProx \cite{li2020fedprox}, FedAdam, an explicit weighted
  adapter (for external FAIR-WEIGHTS-H weights), Byzantine-robust aggregation
  (Krum \cite{blanchard2017krum}, multi-Krum, trimmed mean, median) with a
  Byzantine detector, a
  pathology-aware adapter, a secure aggregator \cite{bonawitz2017secure}, and a
  FAIR-WEIGHTS-H aggregator.
\item \textbf{Privacy:} differential-privacy engines (including an Opacus-backed
  DP-SGD engine), privacy budget tracking and accounting, and noise calibration.
\item \textbf{Monitoring and anomaly detection:} coordinator monitoring plus
  Byzantine detection and FAIR-WEIGHTS-H per-parameter-group risk and reason
  codes.
\item \textbf{Secure communication and aggregation:} TLS-based gRPC with
  certificate validation and mutual auth, and a TenSEAL-backed secure
  aggregation protocol that rejects plaintext in the aggregate-only path.
\item \textbf{Heterogeneity and systems concerns:} FedProx for non-IID
  proximal terms, async training with staleness weighting and adaptive
  timeouts, fault tolerance with checkpoint management, network monitoring, and
  reconnection handling, compression (quantization, sparsification), and
  client resource awareness.
\end{itemize}

\subsection{Execution validation}

PathologyFL has integration and property tests that exercise five-client
training, convergence, DP budget enforcement, Byzantine attacks, client dropout
and checkpoint recovery, weighted aggregation, secure-aggregation homomorphism,
and compression. A PCam-derived federated smoke run uses real PCam patches split
into five simulated sites and four weighting strategies.

\subsection{Evidence status}

The framework is \emph{implemented research infrastructure} with validated
smoke and integration behavior. It is \emph{not} a validated multi-center
deployment system. Key distinctions are enforced:

\begin{itemize}
\item implemented infrastructure and tested execution are separate from
  simulated pathology benchmarks;
\item no real multi-center clinical deployment exists;
\item PCam smoke results are plumbing validation, not performance evidence;
\item centralized frozen-feature CAMELYON17 studies are not described as full
  federated-learning validation.
\end{itemize}

\noindent Two caveats are documented. The FAIR-WEIGHTS-H aggregator integration
(\texttt{pathoalign\_fair.py}) has no dedicated test file, and the
\emph{e2e} federated tests contain dangling imports of modules that do not
exist, so those tests cannot execute as written. Differential-privacy and
secure-aggregation guarantees are gated on optional libraries and are not
independently audited.
